\section{Discussion}
\label{sec:discussion}

\paragraph{Why did mutation operators apply across diverse domains?}
At first glance, 3D mesh decoding, game rendering, and HTTP web serving appear to be unrelated domains.
However, their vulnerability profiles share common structural mechanisms.
All three domains process hierarchical data containing lengths, offsets, and binary buffers.
An operator injecting extreme integer boundaries into glTF buffer views is equally effective at triggering integer overflows in \raylib audio calculations or \drogon \texttt{Content-Length} headers.
Our empirical results indicate that boundary value and chunk mutators find utility across domain boundaries when paired with an adaptive bandit that learns arm weights online.

\paragraph{Do headless in-memory shims miss real bugs?}
By abstracting operating system sockets and graphics displays, headless shims bypass kernel TCP stacks, TLS handshakes, connection reuse state machines, and hardware timing.
Consequently, in-memory shims cannot detect kernel-level socket race conditions, TLS negotiation bugs, or hardware-specific GPU driver faults.
However, user-space memory corruptions, buffer overflows, and parser desynchronizations reside in application code.
Testing in memory enables executing hundreds of millions of iterations that are computationally impractical over loopback sockets.

\paragraph{Why are \texttt{noexcept} filesystem operations a severe threat to web server reliability?}
C++ developers frequently mark file management methods \texttt{noexcept} to satisfy interface contracts or signal intended exception safety.
However, standard filesystem functions throw \texttt{std::filesystem\_error} upon encountering unexpected runtime conditions, such as filenames exceeding length limits.
When an exception escapes a \texttt{noexcept} boundary, the runtime aborts the process via \texttt{std::terminate()}.
Because web servers run multiple worker threads within a single address space, a single malformed upload crashes the entire service.
Developers must consistently use non-throwing \texttt{std::error\_code} overloads in all paths processing untrusted inputs.

\paragraph{Could static analysis have detected defect D1?}
Clang-Tidy includes a check (\texttt{bugprone-exception-escape}) designed to flag functions declared \texttt{noexcept} that invoke throwing operations.
However, in large codebases with nested filesystem wrappers and templates, static checkers generate high false-alarm rates or require whole-program call graphs.
In contrast, coverage-guided fuzzing produced an indisputable crashing reproducer with zero false positives.

We leave the automated synthesis of headless mock shims from interface definitions for future work.

\subsection{Threats to Validity}
\label{sec:threats}

\paragraph{Internal Validity.}
Fuzz testing is non-deterministic.
To mitigate this threat, we executed 10 independent trials on \fastgltf reporting Mann-Whitney U tests and Vargha-Delaney $\hat{A}_{12}$ effect sizes.
We ran multi-worker campaigns totaling 804 million inputs on \raylib and \drogon, and verified all defects with standalone C++ reproducers.

\paragraph{External Validity.}
Our evaluation examines three open-source C and C++ projects: \fastgltf, \raylib, and \drogon, which may not represent all C/C++ software.
To mitigate this threat, we selected mature systems spanning three distinct architectural domains: 3D binary file parsing, interactive multimedia rendering, and high-concurrency network servers.
However, our conclusions may not directly generalize to bare-metal embedded firmware or safety-critical operating systems.

\paragraph{Construct Validity.}
A potential threat is whether branch edge coverage reflects software security testing effectiveness.
To mitigate this threat, we paired coverage feedback with AddressSanitizer and UndefinedBehaviorSanitizer instrumentation.
Furthermore, we evaluated bug-finding effectiveness by reporting candidate defects to upstream maintainers and confirming their security impact through merged pull requests and official bug trackers.
